\subsection{All row permanent quadratic minima and a polynomial ten-plane curve}
\begin{finding}{P-20}{row permanent quadratic minima and a ten-plane degeneration}
Let $k$ be any algebraically closed field, let $n\ge2$, and put
\[
 U^{\mathrm{row,per}}_{n,2}
   =\langle e_i\otimes e_j+e_j\otimes e_i:0\le i<j<n\rangle.
\]
For every $1\le j\le\binom n2$, the exact Grassmannian minimum is
\[
 \mu_j(U^{\mathrm{row,per}}_{n,2})=j+1.
\]
In addition, the polynomial permanent coordinate ten-plane with rectangle subset
\[
 S_{10}=\{0,1,2,3,4,6,7,8,9,10\}
\]
has residual border rank at most $18$.  Consequently
\[
 \mu_{10}(U^{\mathrm{per}}_{4,2})\le18
 \quad\text{in every characteristic},
 \qquad
 17\le\mu_{10}(U^{\mathrm{per}}_{4,2})\le18
 \quad\text{in characteristic zero}.
\]
This does not decide between $17$ and $18$.
\end{finding}
\begin{proof}
For a vector $v$ put $q(v)=v\otimes v$ and write
$s_{ij}=e_i\otimes e_j+e_j\otimes e_i$.
We first prove the general row statement.
Every matrix in $U^{\mathrm{row,per}}_{n,2}$ is symmetric with zero
diagonal.  A nonzero rank-one symmetric matrix has the form
$\lambda vv^{\mathsf t}$, since symmetry makes its two rank-one
factors proportional.  Its diagonal cannot vanish over a field unless
the matrix is zero.  Hence every nonzero matrix in every $j$-plane
$W\subset U^{\mathrm{row,per}}_{n,2}$ has rank at least two.
The residual inclusion tensor has first-factor dimension $j$, so
R-01 with one retained dimension gives
$\underline R(P_W)\ge j-1+2=j+1$.  This argument is valid also in
characteristic two.

For the matching upper bound, choose a graph $G$ of $j$ edges containing
a full star centered at vertex $0$ on its used vertices
$0,\ldots,v-1$.  Such a graph exists for every allowed $j$:
if $j<n-1$, take $v=j+1$ and the star itself; otherwise take $v=n$,
all $n-1$ star edges, and any $j-(n-1)$ of the other edges.
Give its edges independent first-factor labels $\gamma_{ij}$ and let
$P_G=\sum_{ij\in G}\gamma_{ij}\otimes s_{ij}$.
For a star edge use
\[
 F_i(t)=\frac{q(e_0+t e_i)-q(e_0)}t=s_{0i}+tq(e_i),
\]
and for a nonstar edge $ij$ use
\[
 F_{ij}(t)=\frac{q(e_0+t(e_i+e_j))-q(e_0+t e_i)
       -q(e_0+t e_j)+q(e_0)}{t^2}=s_{ij}.
\]
All the star points needed by the second formula are present because
$G$ contains the full star.  After collecting coefficients, the curve
\[
 P_G(t)=\sum_{0i\in G}\gamma_{0i}\otimes F_i(t)
          +\sum_{ij\in G,\ i,j>0}\gamma_{ij}\otimes F_{ij}(t)
       =P_G+t\sum_{i=1}^{v-1}\gamma_{0i}\otimes q(e_i)
\]
uses the base point $q(e_0)$, $v-1$ star points, and one point for
each of the $j-(v-1)$ nonstar edges: exactly $j+1$ pure summands.
The identities hold over $\mathbb Z[t,t^{-1}]$ and use no inverse of
an integer, so the degeneration works in every characteristic.
Its regular tensor-valued limit proves
$\underline R(P_G)\le j+1$ and finishes the exact row minimum.

We now specialize this construction to the five-edge graph
$\{01,02,03,12,13\}$, whose row tensor is
\[
 G=\gamma_{01}\otimes s_{01}
   +\gamma_{02}\otimes s_{02}
   +\gamma_{03}\otimes s_{03}
   +\gamma_{12}\otimes s_{12}
   +\gamma_{13}\otimes s_{13}
 \in k^5\otimes k^4\otimes k^4.
\]
For $t\ne0$ define
\begin{align*}
 F_i(t)&=\frac{q(e_0+t e_i)-q(e_0)}{t}
       =s_{0i}+tq(e_i),\qquad i=1,2,3,\\
 F_{ij}(t)&=\frac{q(e_0+t(e_i+e_j))-q(e_0+t e_i)
                    -q(e_0+t e_j)+q(e_0)}{t^2}
       =s_{ij},\qquad (i,j)=(1,2),(1,3).
\end{align*}
These are identities over $\mathbb Z[t,t^{-1}]$ and involve no division
by an integer.  Put
\[
 G(t)=\sum_{i=1}^3\gamma_{0i}\otimes F_i(t)
       +\gamma_{12}\otimes F_{12}(t)
       +\gamma_{13}\otimes F_{13}(t).
\]
After collecting the first-factor coefficients, it uses exactly the
six displayed square vectors
\[
 q(e_0),\ q(e_0+t e_1),\ q(e_0+t e_2),\ q(e_0+t e_3),
 q(e_0+t(e_1+e_2)),\ q(e_0+t(e_1+e_3)).
\]
Each has one arbitrary first-factor coefficient and therefore gives
one pure three-factor summand.  Thus $\operatorname R(G(t))\le6$, and
\[
 G(t)=G+t\sum_{i=1}^3\gamma_{0i}\otimes q(e_i)
 \longrightarrow G.
\]
The sum has a regular limit even though the individual first-factor
coefficients may have poles.

The general row result just proved shows that this factor has border
rank exactly six.

The two-edge star tensor
\[
 H=\alpha_1\otimes s_{01}+\alpha_2\otimes s_{02}
       \in k^2\otimes k^3\otimes k^3
\]
has the three-term curve
\[
 H(t)=\alpha_1\otimes F_1(t)+\alpha_2\otimes F_2(t)
      =H+t\sum_{i=1}^2\alpha_i\otimes q(e_i).
\]
It uses $q(e_0),q(e_0+t e_1),q(e_0+t e_2)$, so its border rank is at
most three, and its support flattening of rank three gives equality
(as also proved in P-03).

Take the Kronecker product of these two explicit curves, grouping
corresponding factors.  For every $t\ne0$, $H(t)\boxtimes G(t)$ is a
sum of at most $3\cdot6=18$ pure tensors in
$k^{10}\otimes k^{12}\otimes k^{12}$, and its regular limit is
$H\boxtimes G$.  Identify $e_r\otimes e_c$ in the last two factors
with the matrix-entry variable $x_{rc}$.  Expanding
$s_{0i}\boxtimes s_{cd}$ gives exactly the four ordered coefficients
of the permanent rectangle with row pair $(0,i)$ and column pair
$(c,d)$.  Here $i=1,2$ and
$(c,d)\in\{01,02,03,12,13\}$.  Their rectangle indices are precisely
$S_{10}$, and the ten first-factor labels are independent.
Thus $H\boxtimes G$ is the required residual inclusion tensor.
This proves the upper bound over every algebraically closed field.

In characteristic zero, R-10 and the retained-dimension propagation
R-07 give $\mu_{10}\ge10+7=17$, proving the stated interval.
The exact script \texttt{check\_per10\_upper18.py} verifies both
integral Laurent identities and every coefficient of the Kronecker
identification; its output is
\texttt{per10\_upper18\_verification.json}.
\end{proof}

\paragraph{Dictionary and remaining gate.}
The constructions are morphisms from the punctured affine line over
$\mathbb Z$ to the respective secant parameterizations whose tensor-valued
sum extends over zero.  Its specialization lies in the secant closure
of the Segre product.  The first and mixed second differences are
integral identities in tensor powers of free modules, so base change
preserves the construction in every characteristic.  The only
characteristic-zero ingredient in the interval is the lower bound
already proved in R-10.  In the common scheme and locally-free-sheaf
dictionary this gives a concrete degeneration, not merely an equality
of $K_0$ classes.  The unresolved ten-plane gate is now whether the
uniform minimum is $17$ or $18$; the curve does not provide a
seventeen-term decomposition.
